\section{Methodology: evidence selection under a fixed budget}
\label{sec:methods}
% src: structure of the earlier version (paper/sections 03_benchmark, 05_frozen); numbers carried over with their original sources
% src: work/datasets_v2/DESIGN.md (tiers, schema, authoring pipeline); questions/*.jsonl (counts); paper/analysis/revision_stats.md "Families" (labels 47, constituents 33, components 20, blocks 29)
% src: RAG-Paper-upgrades/ragbench: retrieval/hybrid.py, config.yaml (chunking), approaches/hybrid.py (rerank), approaches/lex5b/{approach,views}.py, approaches/lex8/fusion.py; work/datasets_v2/setr/SETR.md; test_freeze/PROTOCOL.md sec. 2 and 9; lex5b/LEX5B.md
% src (number check 2026-11-11): 48 of 54 = work/datasets_v2/lex8/LEX8.md "Diagnosis" (also lex9/LEX9.md); 32 of 156 (31 good, 1 fix) = work/TODO.md line 95; 35 of 4,608 unresolved = TEST_RESULTS.md sec. 6 and S5_stats/RESULTS.md sec. 0; pool chars = work/multidoc/pools/<pool>.json (fa 11.64M, en 26.42M)
% Shortened 2026-11-18 (length cap): implementation detail, the pools and settings tables and the model cross-check cut from the body; earlier text in backups/pre_cut_2026-11-18/sections/.

A retrieval-augmented generation (RAG) system answers a question in three stages \citep{lewis2020rag,gao2023ragsurvey}. A retrieval stage searches a document collection for passages that may contain the answer, a selection stage chooses which of them to show, and a reading stage lets a language model, called the reader, write the answer from the selected text. This study varies the first two stages and holds the third fixed. Figure~\ref{fig:pipelines} shows the four pipelines that are compared. Each splits the documents into chunks and ranks them with BM25, dense vectors or both, may reorder or filter the ranked chunks with a cross-encoder or a language model, adds the chosen chunks to an evidence pack of limited length (a step called packing), and passes the pack to the reader. A judge then scores the answer against a rubric.
% src: code-derived schematic, no estimated quantities. RAG-Paper-upgrades/ragbench: retrieval/hybrid.py,
%      approaches/hybrid.py (RRF depth 40; rerank: 50 fused candidates rescored, top 40 packed), approaches/lex5b/approach.py
%      (depth 40, top 24 shown, view 'literal', 3 s deadline, fallback = fusion order), approaches/lex5/bundles.py (exact
%      bundle packing + fill), approaches/lex5b/views.py, approaches/setr/approach.py (top 20, PASSAGE_CAP 1500, 8 s,
%      fallback = HYBRID order); test_freeze/PROTOCOL.md sec. 2 and 9 (readers, judge, human review GLM 6,000, 30 q).
\begin{figure*}[tp]
\centering
\begin{tikzpicture}[
  x=1cm, y=1cm,
  box/.style={stage=#1, text width=2.25cm, minimum height=1.3cm},
  neutral/.style={draw=cAxis!70, line width=0.8pt, rounded corners=4pt, fill=black!4, align=center, font=\footnotesize,
    inner sep=3pt, text width=2.25cm, minimum height=1.3cm, text=cInk},
  ghost/.style={draw=cAxis!45, line width=0.8pt, dashed, rounded corners=4pt, align=center, font=\footnotesize,
    inner sep=3pt, text width=2.25cm, minimum height=1.3cm, text=cAxis},
  head/.style={font=\small\bfseries, text=cInk, anchor=south, text width=2.6cm, align=center},
  tag/.style={draw=#1, fill=white, line width=0.8pt, rounded corners=2pt, font=\scriptsize\bfseries, text=#1,
    inner sep=1.6pt},
]
\def\cA{3.05}\def\cB{5.85}\def\cC{8.65}\def\cD{11.50}
% ---------------------------------------------------------------- column headers
\node[head] at (\cA,1.12) {Candidates};
\node[head] at (\cB,1.12) {System-specific step};
\node[head] at (\cC,1.12) {Evidence pack ($\le B$)};
\node[head] at (\cD,1.12) {Reader and scoring};
\draw[cGrid, line width=0.6pt] (0.1,1.06) -- (12.70,1.06);
% ---------------------------------------------------------------- lane labels
\foreach \y/\c/\n/\s in {0/cHyb/\hybrid/no extra call, -2.9/cHR/\hr/cross-encoder, -5.8/cLex/\lexb/1 LLM call, -8.7/cSetR/\setr/1 LLM call}{
  \node[lanelabel=\c] at (1.8,\y+0.14) {\n};
  \node[font=\footnotesize, text=cAxis, anchor=east] at (1.8,\y-0.26) {\s};
}
% ---------------------------------------------------------------- candidates
\node[neutral] (r1) at (\cA,0)     {BM25 top 40 + dense top 40, rank fusion (RRF)};
\node[neutral] (r2) at (\cA,-2.9) {same retrieval; up to 50 fused candidates};
\node[box=cLex] (r3) at (\cA,-5.8) {BM25 + dense + learned sparse; score fusion, 40 evidence units};
\node[neutral] (r4) at (\cA,-8.7) {\hybrid{} candidates: top 20 chunks, shown whole};
% ---------------------------------------------------------------- system-specific step
\node[ghost] (s1) at (\cB,0) {no selection step};
\node[box=cHR] (s2) at (\cB,-2.9) {cross-encoder rescores the 50; keeps the top 40};
\node[box=cLex] (s3) at (\cB,-5.8) {reads excerpts of the top 24 of the 40 evidence units; names bundles};
\node[box=cSetR] (s4) at (\cB,-8.7) {SetR prompt with chain of thought selects a set};
\node[tag=cHR]   at (s2.north west) [xshift=0.42cm, anchor=south, yshift=-0.5pt] {CE};
\node[tag=cLex]  at (s3.north west) [xshift=0.42cm, anchor=south, yshift=-0.5pt] {LLM};
\node[tag=cSetR] at (s4.north west) [xshift=0.42cm, anchor=south, yshift=-0.5pt] {LLM};
\node[tag=cLex]  at (s3.north east) [xshift=-0.62cm, anchor=south, yshift=-0.5pt] {3\,s limit};
\node[tag=cSetR] at (s4.north east) [xshift=-0.62cm, anchor=south, yshift=-0.5pt] {8\,s limit};
% ---------------------------------------------------------------- packing
\node[neutral] (p1) at (\cC,0)     {ranked chunks, fused order; last chunk clipped};
\node[neutral] (p2) at (\cC,-2.9) {ranked chunks, reranked order; last chunk clipped};
\node[box=cLex] (p3) at (\cC,-5.8) {exact bundle packing of whole units, then fill; no clipping};
\node[neutral] (p4) at (\cC,-8.7) {selected set first, then \hybrid{} fill; last chunk clipped};
% ---------------------------------------------------------------- flows
\foreach \i/\c in {1/cHyb,2/cHR,3/cLex,4/cSetR}{
  \draw[flow=\c] (r\i) -- (s\i);
  \draw[flow=\c] (s\i) -- (p\i);
}
% fallbacks: a missed time limit or an unusable reply skips the selection
\draw[fallback=cLex] (r3.south) -- ++(0,-0.17) -| node[pos=0.25, fill=white, inner sep=1pt, font=\footnotesize, text=cLex] {fallback: fusion order} (p3.south);
\draw[fallback=cSetR] (r4.south) -- ++(0,-0.17) -| node[pos=0.25, fill=white, inner sep=1pt, font=\footnotesize, text=cSetR] {fallback: \hybrid{} order} (p4.south);
% ---------------------------------------------------------------- shared reader and scoring
\node[shared, text width=2.2cm, minimum height=2.6cm] (reader) at (\cD,-1.45)
  {\textbf{Reader}\\[1pt] one answer prompt; GLM-5.3-flash (primary); 5 more readers answer the same packs};
\node[shared, text width=2.2cm, minimum height=2.6cm] (judge) at (\cD,-7.25)
  {\textbf{Scoring}\\[1pt] strict LLM judge on every answer; human review: GLM, 6,000 chars, 30 selected questions};
\coordinate (bus) at (10.12,0);
\foreach \i in {1,2,3,4}{\draw[line width=1.1pt, cAxis] (p\i.east) -- (p\i.east -| bus);}
\draw[line width=1.1pt, cAxis] (p1.east -| bus) -- (p4.east -| bus);
\draw[flow] (bus |- reader.west) -- (reader.west);
\draw[flow] (reader) -- (judge);
\end{tikzpicture}
\caption{The four pipelines. All share chunking, an answer prompt and a maximum evidence budget $B$, and coloured boxes mark steps specific to one pipeline. The two language-model selectors fall back on a timeout or an unusable reply. Five more readers answer the packs of the primary run. The selector of the proposed pipeline has a deadline of 3 seconds and that of the set selector of 8 seconds.}
\label{fig:pipelines}
\end{figure*}


\subsection{Problem formulation}
\label{sec:formulation}
The setting is a closed search over fixed document pools. A pool $P$ is a set of documents in one language and one genre, and $\mathcal{C}(P)$ denotes the set of all chunks of the pool. A question $q$ is asked in the language of its pool and is answered against that pool alone. Each question has a tier $t(q)\in\{1,2,3\}$ that states how widely the needed evidence is spread (Sect.~\ref{sec:benchmark}), and an annotation records one or more evidence bundles $\mathcal{G}(q)=\{G_1,\dots,G_m\}$, where a bundle is a set of passages that together suffice to answer the question.

A pipeline $\pi$ maps a question to an evidence pack $E=\pi(q,P,B)$, a list of chunks from $\mathcal{C}(P)$. The evidence budget $B$ limits the length of the pack in rendered characters, so that a pack is feasible if
\begin{equation}
\sum_{c\in E}\ell(c)\;\le\;B,
\label{eq:budget}
\end{equation}
where $\ell(c)$ is the length of chunk $c$ with its header. Equal characters are not equal tokens: the same text costs more tokens in Persian than in English. A reader writes an answer from the question and the pack alone, and a judge returns a verdict of 1 for a correct answer and 0 otherwise. Accuracy is reported as the number of correct answers over a question set. Because accuracy mixes the quality of the pack with that of the reader, a second quantity measures the pack alone. A pack is complete for a question if it contains every passage of at least one evidence bundle:
\begin{equation}
\mathrm{comp}(q,E)\;=\;\mathbb{1}\bigl[\exists\,G\in\mathcal{G}(q):\;G\subseteq E\bigr].
\label{eq:complete}
\end{equation}
The share of complete packs shows how often a pipeline gives the reader enough to answer. It also separates two kinds of loss: evidence can be lost before a candidate list exists, when retrieval misses it, or afterwards, when packing drops it. The pipelines are compared at the same budget, and Sect.~\ref{sec:evidence} varies the budget itself.

\subsection{Benchmark}
\label{sec:benchmark}
The benchmark pairs natural Persian and English documents with questions of three difficulty tiers, and each question is answered against its whole pool. It complements existing Persian and multilingual resources \citep{darvishi2023pquad,khashabi2021parsinlu,zinvandi2025famteb,zhang2023miracl}, and its tiers resemble the task types of multi-hop and cross-document question sets \citep{yang2018hotpotqa,tang2024multihoprag}. The documents come from three genres in both languages, namely company reports, historical or expository non-fiction, and statutes, and only born-digital text or verified transcriptions were used. Table~\ref{tab:tiers} defines the tiers. The tier of a question is the least widely spread evidence that suffices to answer it: one passage for T1, several regions of one document for T2, and at least two documents for T3. The six search pools hold approximately 12 million characters of Persian and 26 million of English text.
% src: work/datasets_v2/DESIGN.md (tier definitions, integration operations); questions/*_{dev,test}.jsonl (counts per tier,
%      language and domain); examples verbatim from questions/law_dev.jsonl V2-dev-LAW-en-T1-003,
%      rpt_dev.jsonl V2-dev-RPT-en-T2-003, rpt_dev.jsonl V2-dev-RPT-en-T3-001, rpt_dev.jsonl V2-dev-RPT-en-AB-001 (dev questions)
\begin{table*}[tp]
\caption{Definition and examples of the three question tiers}
\label{tab:tiers}
\centering
\small
\renewcommand{\arraystretch}{1.15}
\begin{tabularx}{\textwidth}{@{}>{\raggedright\arraybackslash}p{2.3cm}>{\raggedright\arraybackslash}p{4.4cm}>{\raggedright\arraybackslash}Xcc@{}}
\toprule
Tier & Simplest sufficient evidence & Example (dev question) & Dev & Test \\
\midrule
\textbf{T1}\newline local fact & one coherent passage, list or table region gives the answer and the labels that bind it to its subject and constraints & \emph{Whose consent is needed before a prosecution under the Bribery Act 2010 can be brought in England and Wales?} (LAW-en) & 18 & 30 \\
\addlinespace
\textbf{T2}\newline within-document integration & one document suffices, but no single region does; definition bridge, rule plus exception, comparison or derivation, bounded synthesis & \emph{Wise's FY2024 annual report gives its total customer balances and, in the Wise Business section, the business deposits it holds. Roughly what share of total customer balances were business deposits?} (RPT-en, derivation) & 18 & 30 \\
\addlinespace
\textbf{T3}\newline cross-document integration & every sufficient bundle spans at least two documents; cross-time or cross-version comparison, cross-subject comparison, cross-document bridge, aggregation & \emph{Turkcell and Vodafone Turkey both had their 2G licences extended to 30 April 2029. What fee did each operator agree for the extension, and how did the way they paid it differ?} (RPT-en, cross-subject) & 18 & 30 \\
\addlinespace
Unanswerable & no sufficient evidence in the pool (absence audit) & \emph{How many members did Zalando Plus have at the end of 2024?} (RPT-en) & 6 & 6 \\
\bottomrule
\end{tabularx}
\end{table*}


% src: tables/pools.tex (Persian pools 4.76+5.70+1.18 = 11.64M characters, English 16.64+6.26+3.52 = 26.42M); 45 answerable test questions per language (tab:tiers); within-language detectable difference about 8 of 45: S14_bilingual/RESULTS.md sec. 2.3 (8.3 of 45 at discordance 0.2)
The split sizes, 60 questions for development and 96 for the test, were set by the number of questions that could be reviewed by hand. Each split adds six unanswerable questions, one per pool, for which a correct answer is an abstention. The primary reader abstained correctly on all six unanswerable test questions in every configuration. % src: R10: TEST_RESULTS.md sec. 2 (AB /6 = 6 for GLM-5.3-flash in all eight rows)
Of the 96 test questions, 90 are answerable, and 45 of them are in each language. Many test questions share source material and are not independent. The 90 answerable questions fall into 20 groups that share a source, and confidence intervals resample whole groups. Each question carries the atomic requirements that an answer must meet, the evidence passages with exact quotes and offsets into the frozen text, and the evidence bundles. A pipeline sees only the question, its language and its pool, plus the expected answer format (a single fact, a list, a comparison, a series or an explanation), which is identical for every pipeline and reader.

The two languages are not translations of each other. They use different documents, so a difference between the halves mixes the effect of the language with the effect of the sources. The design therefore supports only properties of a language or its tools that can be measured directly, such as the token cost of a character, findings that appear in both halves, and designs in which every question is its own control, such as asking every question in the other language. A paired comparison within one half detects only a difference of approximately 8 answers out of 45, so any other split by language is descriptive.

The questions were built in six steps, four automated and two by people. Claude Opus 5.5 agents wrote a specification for each question and drafted the question from it, in line with earlier use of language models to generate queries \citep{bonifacio2022inpars,dai2023promptagator}. A mechanical verifier and a separate source-review agent then checked the drafts. A native Persian speaker reviewed the Persian wording, and the author spot-checked 32 of the 156 questions, of which 31 were accepted and 1 was fixed. Development and test share no evidence document, source block or information need, but they share the searchable pools. The schema and the construction steps are in \texttt{DESIGN.md} and \texttt{SCHEMA.json} of the benchmark dataset, and the construction prompts are in the code repository (\texttt{code/corpus\_build}).

\subsection{Retrieval and selection pipelines}
\label{sec:systems}
All pipelines share the same chunking (at most 1,000 characters with an overlap of 200), answer prompts and readers, so whole pipelines are compared and a difference is not credited to a single component. Model versions are in \texttt{protocol/reader\_pins.json} and \texttt{protocol/judge\_pin.json} of the results dataset, and the other parameter values are in \texttt{code/ragbench/config.yaml} of the code repository.

The first pipeline, \hybrid{}, combines a lexical and a dense retrieval channel. The lexical channel is BM25, which scores a chunk by the weighted overlap of its terms with the query \citep{robertson2009bm25}. The dense channel ranks chunks by the similarity of their vectors to that of the question, using bge-m3, a multilingual encoder that also produces learned-sparse term weights \citep{chen2024bgem3,karpukhin2020dpr}. \hybrid{} merges the two ranked lists by reciprocal rank fusion (RRF) \citep{cormack2009rrf}, which scores a chunk by the sum of $1/(60+\mathrm{rank})$ over the two lists. It adds the chunks to the pack in fused order until the budget is reached, and it clips the last chunk if enough of it fits.

The second pipeline, \hr{}, is \hybrid{} with a reranker. A cross-encoder reads the question and one chunk together and outputs a relevance score \citep{nogueira2019monobert}, and the reranker used here is bge-reranker-v2-m3 on a local GPU. \hr{} packs the chunks in reranked order with the same clipping rule and needs no tokens from a language model.

The third pipeline, \lexb{}, is the proposed pipeline, and its name is a label that does not mean that the method is purely lexical. It has four stages. The first stage merges BM25, dense and learned-sparse scores into one ranking. Each channel score is min-max normalised over the union of the three top-40 lists and clipped to the interval from 0 to 1, which gives $z_j(c)$, and the fusion score is
% src: ragbench/approaches/lex8/fusion.py Normalizer.z (min-max over the union of the top-40 lists, clipped to [0,1], defined for every chunk from the full score array) and maxmean
\begin{equation}
s_{\mathrm{fuse}}(c)\;=\;\tfrac12\max_{j}z_j(c)\;+\;\tfrac12\cdot\frac{1}{3}\sum_{j=1}^{3}z_j(c).
\label{eq:fuse}
\end{equation}
Unlike rank fusion, this score uses the actual channel scores, so it responds more strongly when one signal is much better than the others \citep{bruch2023fusion}. The second stage turns the top chunks into evidence units by extending sentences that are cut at a chunk edge. In the third stage, one call to a language model, the selector, reads short excerpts of the top 24 units and splits the question into at most eight parts. For each part it names a bundle of units that together support it, with at most one alternative bundle. In the fourth stage, let $S$ be a choice of at most one listed bundle for each part, let $\mathrm{cov}(S)$ be the number of parts that $S$ covers, and let $\mathrm{alt}(S)$ be the number of covered parts that rest on an alternative bundle. The packer solves
\begin{equation}
S^{\star}\;=\;\operatorname*{arg\,max}_{S:\;\mathrm{len}(S)\le B}\;\bigl(\mathrm{cov}(S),\;-\mathrm{alt}(S),\;-\mathrm{len}(S)\bigr),
\label{eq:pack}
\end{equation}
where the tuple is compared in lexicographic order and $\mathrm{len}(S)$ is the rendered length of the chosen units. The remaining budget is filled with the next whole units of the fused ranking. If the call fails, misses its 3-second deadline or breaks the required reply format, the pipeline falls back to the fused ranking.

The fourth pipeline, \setr{}, adapts the published SetR selector \citep{lee2025setr}, which replaces a ranked list by a selected set. One chain-of-thought call uses the prompt of SetR to pick a set from the top 20 chunks of \hybrid{}, the selected chunks come first, and the list of \hybrid{} fills the rest of the budget. The letter p stands for prompted, because an untuned general model receives the SetR prompt in place of the fine-tuned selector of the original, and fill stands for the filling of the budget. The results therefore do not describe the published system.

Several variants around \hybrid{} and \hr{} change one component at a time. They vary the retrieval channel (BM25 or dense retrieval alone), query rewriting (a hypothetical answer passage \citep{gao2023hyde} or three paraphrases \citep{rackauckas2024ragfusion}), chunk size (500 and 2,000 characters), and the embedding model (multilingual-e5-large \citep{wang2024me5}). Further variants remove the selection call of \lexb{}, replace it by a rule without a language model, apply a listwise reranker \citep{sun2023rankgpt,zhuang2024setwise}, let a model search the pool up to five times before it names passages \citep{yao2023react,jin2025searchr1}, pass whole documents to the reader \citep{li2024selfroute,xu2024retrieval}, or route each question to a cheaper or a more expensive pipeline. Every variant uses the same chunks, prompts, reader and judge as the four main pipelines.

\subsection{Readers and judging}
\label{sec:readers}
% src: judge 2 on the 8 frozen cells = S1_evaluator/judge2_results.md sec. 1; on 7 post hoc variants (FUSION_6000, SETR_NOFILL, LISTWISE, ITER_BM25, GREP_AGENT, GOLD, CLOSEDBOOK) = S8_factorial/RESULTS.md l.66-69; 35 of 4,608 = S5_stats/RESULTS.md l.11; up to five searches = S3_agentic/RESULTS.md l.132
One reader, GLM-5.3-flash, answers every question and also makes both selection calls. Five readers from other model families (gpt-oss-20b, Qwen3-235B-A22B-2507, Llama-3.3-70B-Instruct, Gemma-3-27B-it and DeepSeek-V4-flash) answer the identical evidence packs. Because they see the same packs, they add breadth but are not independent replications. All readers are open-weight models that were called through an API, each fixed to one provider and model version.

Answers are scored by an LLM judge, a language model that grades an answer against a reference and a rubric \citep{zheng2023judge}. Such a judge is cheap and scales to many answers, but its verdicts can vary between runs and can differ from those of a person \citep{faggioli2023perspectives,thakur2025judging,haldar2025roulette}. The first judge, gpt-oss-120b, applies a strict rubric: an answer counts as correct only if it meets every required point and contains no contradiction, wrong value, wrong subject or unsupported added claim. A cell is one answer of one reader to one question in one configuration of pipeline and budget. The judge sometimes returned a reply that could not be parsed after three attempts. These 35 of the 4,608 cells (96 questions, eight configurations and six readers) are unresolved, and scores are reported as ranges that cover every possible verdict for them. A second judge from a different model family, Mistral Large, graded the answers of the primary reader in the eight frozen configurations and in seven post hoc variants: \lexb{} without selection, the gold-evidence run, the closed-book run, \setr{} without the fill, the listwise reranker, the BM25 agent and the regular-expression agent. % src: R10: S8_factorial/RESULTS.md l.66-69 (FUSION_6000, SETR_NOFILL, LISTWISE, ITER_BM25, GREP_AGENT, GOLD, CLOSEDBOOK)
The judge prompt is in \texttt{code/ragbench/eval/judge.py} of the code repository.

\subsection{Experimental setup}
\label{sec:protocol}
% src: 15,500 = ragbench/approaches/lex10/approach.py docstring (largest 500-step pack with predicted answer-call tokens <= 2.40x HYBRID; lex10/tokbudget.py); +5 of 90 from +3/54 and 1.5x -> 2.5x = PROTOCOL.md sec. 5 l.73-78 and lex5b/LEX5B.md l.3-4
All method development used the development questions, and Claude Opus 5.5 agents implemented and ran the methods without reading the test files. Before the first test answer, the code, its inputs and the decision rules were frozen, and each test question was run once in every pipeline and budget. The freeze is a commit in the repository of the author and was not registered with a third party, although preregistration is advocated for NLP studies \citep{vanmiltenburg2021prereg}. The test used a tight budget of 6,000 characters and a loose budget of 15,500 characters, the largest pack whose predicted reader prompt stayed within the token limit of the protocol. Two gates and two tests were fixed in advance, all for the primary reader. The accuracy gate required \lexb{} at the tight budget to beat each control, \hybrid{} and \hr{}, by at least 5 of the 90 answerable questions, and the resource gate limited the query-time tokens of \lexb{}. Two superiority tests covered the main comparisons, and every other analysis was descriptive. The most important protocol change relaxed the token limit from 1.5 to 2.5 times the tokens of \hybrid{}, before any \lexb{} result existed. Every change is logged in \texttt{protocol/PROTOCOL.md} of the results dataset.

A human review complements the judge. The author reviewed the answers of the primary reader at 6,000 characters on the 30 questions for which the judge gave the four pipelines different verdicts, with the pipeline hidden. On each of these questions the answers of all four pipelines were reviewed, which gives 120 answers. This review is called the complete review. An earlier review, called the unequal review, also covered questions on which the pipelines disagreed, but not every pipeline on each of them, and Sect.~\ref{sec:stability} explains why that matters. On the reviewed questions the human verdicts replace the judge's, and on all other questions the judge's verdicts stand. The resulting merged scores were built after the judge scores and are therefore post hoc.

% src: 90% / 12-character / same-document presence rule = S7_analysis/RESULTS.md l.21 and s7lib.py MIN_SEG = 12; 30 questions, 120 answers = TEST_RESULTS_HUMAN_SYM.md "Coverage"
The main metric is the number of correct answers out of the 90 answerable questions, and pack quality is the number of questions with a complete pack (Eq.~\ref{eq:complete}). A gold evidence item counts as present in a pack if at least 90\% of its quote appears word for word, in pieces of at least 12 characters, inside passages of its own document. Cost is the number of tokens of all query-time language-model calls per question, including reasoning tokens, because pricing and speed follow tokens. The cross-encoder, the embedding model, index construction and storage are not counted. Latency was measured while other jobs were running on the same machine and is descriptive.
% src: S5_stats/RESULTS.md sec. 1 (estimators written in numpy/scipy; Laplace approximation; Wald SEs from a finite-difference Hessian; statsmodels mixed GLM is variational Bayes only; lme4 not used)
Two pipelines are compared on the same questions with a paired analysis. The test is the two-sided exact sign test on the questions on which only one of the pipelines is correct, which is the exact form of McNemar's test \citep{mcnemar1947}. Confidence intervals are percentile bootstrap intervals from 10,000 resamples of the paired difference, and the cluster version resamples the 20 groups of questions that share a source \citep{efron1993bootstrap}. Multiple comparisons are corrected with the Holm procedure where stated \citep{holm1979}, and all other $p$-values are uncorrected and descriptive. Pooled analyses use a mixed-effects logistic regression of every answer of all six readers, with pipeline and budget as fixed effects and random intercepts for the question and for the evidence pack, because the six readers answered identical packs \citep{bates2015lme4}. It was fitted by maximum likelihood with a Laplace approximation in Python and cross-checked with a variational Bayes fit and with generalised estimating equations in the statsmodels library, which give effects in the same direction. The smallest difference that the design can detect depends on the number of questions $N$ and on the rate $d$ at which two pipelines disagree \citep{card2020power,sakai2016topicset}. A normal approximation to the paired test gives the minimum detectable effect in answers for a two-sided level $\alpha$ and power $1-\beta$:
\begin{equation}
\Delta_{\min}\;\approx\;\bigl(z_{1-\alpha/2}+z_{1-\beta}\bigr)\sqrt{N\,d}.
\label{eq:mde}
\end{equation}
% src: S5_stats/RESULTS.md sec. 2.3 (psi .149/.170; MDE exact 10.4/11.1, normal 10.1/10.8 of 90); S14_bilingual/RESULTS.md sec. 2.3 (n = 45: 6.8 at psi .15, 8.3 at psi .20)
With $N=90$, $d$ between 0.15 and 0.17, $\alpha=.05$ and 80\% power, the value is approximately 11 answers. A half of the benchmark with $N=45$ detects approximately 7 to 8.

% src: 576 packs = S2a_packs/RESULTS.md l.16; families: S13_embed 6, S11_rewrite 7, S10_retrieval 12, S12_chunks 13 paired rows (RESULTS.md); agents 17 = S3_agentic/RESULTS.md l.95; language 14 = S14_bilingual/RESULTS.md l.89
All language models were reached through the OpenRouter service, each through one fixed provider and model version, with temperature 0 and a fixed seed where the provider accepts one. GLM-5.3-flash has no seed parameter, so its answers are not deterministic. Each experiment of Sect.~\ref{sec:results} ran once. The answers of the frozen test are called the saved run, and the repeated runs of Sect.~\ref{sec:stability} send the same packs to the reader again. Only the frozen test is confirmatory, because its hypotheses, prompts and decision rules were fixed before any test answer was seen. All other experiments use the same test questions after their answers were known and are therefore post hoc. Their prompts and settings were fixed on the development set or before the run. Packs rebuilt from the saved traces with the frozen code matched the frozen packs byte for byte (576 packs, which are the six configurations of \hybrid{}, \hr{} and \lexb{} at both budgets, each over 96 questions). Each component experiment compared a declared family of 6 to 13 pairs, and the agent and language analyses used families of 17 and 14. Because the same 90 questions serve all of them, post hoc $p$-values are descriptive \citep{dwork2015reusable}.
